
Neural network model repositories contain large numbers of models whose architectural properties are not always reliably documented. Determining basic characteristics such as network depth typically requires either parsing metadata that may be incomplete or absent, or executing forward passes that require test data and computational resources. This creates a practical obstacle for practitioners who need to verify architectural properties before model deployment or selection.

Prior work on neural network analysis has focused primarily on weight distributions for compression and transfer learning purposes. Weight pruning methods identify redundant parameters through magnitude-based criteria. Quantization research analyzes weight value distributions to optimize bit-width allocation. Transfer learning studies measure weight similarity across models to predict transfer performance. These approaches treat weight distributions as properties to be optimized or transferred, not as indicators of architectural characteristics.

This study examines whether architectural depth can be classified from weight statistics alone. We formulate the problem as binary classification: shallow networks (≤34 layers) versus deep networks (≥50 layers). We test three types of weight-derived features on 20 pretrained convolutional neural networks spanning ResNet, VGG, and DenseNet families. All models were pretrained on ImageNet-1K using standard training protocols available through PyTorch torchvision.

The research questions are: (1) Do weight statistics enable binary depth classification above chance level (50\%)? (2) Are discriminative features architectural (present in network structure) or training-induced (emergent from gradient flow)? (3) Do depth signals persist within single architecture families when controlling for family-specific design patterns?

Our findings include: Test accuracy of 100\% on 4 held-out models for global statistics, layer-wise features, and architectural features. Random initialization testing showing zero performance degradation when using untrained models, indicating architectural rather than training-induced features. Within-family validation accuracy of 100\% for both ResNet and DenseNet families. Feature importance analysis identifying layer count and bottleneck ratio as primary discriminators.
